\section{A Name Which No Man Knoweth}

The risen Christ of the Gospels is, if anything, more stubbornly particular than the crucified one. The Easter narratives could easily have dissolved him into radiance; instead John shows him on a beach at dawn, having made breakfast: ``they saw hot coals lying, and a fish laid thereon, and bread.'' The glorified Lord of the universe has cooked. And the restoration of the apostle who denied him is not conducted by general amnesty but by name, in a thrice-repeated question aimed at one man's specific wound: ``Simon son of John, lovest thou me more than these?''\endnote{John 21:9, 12, 15, Douay--Rheims, checked 24 July 2026. The threefold question answering the threefold denial is a classical reading of the passage's structure, here used as reading rather than asserted as the evangelist's stated intention.} Whatever resurrection means, it does not mean graduation from particularity. It means particularity made permanent.

This is the point at which the objection's deepest premise can finally be named. Behind the preference for the universal over the particular lies a picture of perfection as generality---as if the best love were the widest, the way the best law is the widest. But love does not behave like law. Nobody wants to be loved in general. To be loved as an instance---of humanity, of womankind, of the deserving poor---is precisely not to be loved; the beloved wants to be seen and wanted as irreplaceably \emph{this} one, and knows instantly the difference between a form letter and a letter. The objector's God, who deals only in universals, could never love in the second way. He could only apply. The scandal of particularity, followed to its end, is the claim that God loves the way lovers love and not the way statutes apply---that the universal cause, who alone knows every singular directly, is thereby the only one for whom \emph{everyone} can be a particular, no one an instance, no one an average.

That is why the promise held out at the end of the New Testament to each of the redeemed is not absorption into a general beatitude but a keepsake of unrepeatable address: ``To him that overcometh, I will give the hidden manna, and will give him a white counter, and in the counter, a new name written, which no man knoweth, but he that receiveth it.''\endnote{Apocalypse (Revelation) 2:17, Douay--Rheims, checked 24 July 2026. The verse is quoted for its image of irreducibly personal address; no claim is made here about the disputed identification of the ``white counter'' (\emph{calculus}, stone) in the history of exegesis.} A name no one else can know is the exact opposite of an abstraction. It is particularity perfected: identity so fully affirmed that it becomes a secret between two.

The argument of this article can now be stated in the order in which it was earned. Nothing exists except particulars, so a God who acts must act among them. Abstraction is the finite knower's economy, not the infinite knower's grandeur, so a God of generalities-only would be less than the classical God, not more. Election, on its own documents, is the choice of one for the sake of all. The Incarnation could not have been general without being nothing, and its one assumed body is the ground on which the Son is united in some fashion with every man. The sacraments prolong the same style, because a wound that is concrete is healed concretely. And the end of the whole economy is not the dissolving of persons into the universal but the white stone: each one named, none interchangeable. The scandal of particularity is real, and it is not a flaw to be explained away. It is the signature. He was always going to be somebody's neighbor, from some town, with a known mother; and the question he asked there is still in force, still addressed---like everything else in this economy---to one person at a time.
